\begin{tabular}{@{}lllr@{}}
\toprule
Question & Quantity & Estimate (95\% interval) & Splits \\
\midrule
\multicolumn{4}{@{}l}{\textit{(a) Are the residuals heavy-tailed? (Q1)}}\\
 & GT $-$ GG, log predictive density per visit & $+0.1115$ ($+0.0606$, $+0.1676$) & 7/7 \\
 & GT $-$ GG, class-label metric              & $+0.0382$ ($+0.0205$, $+0.0549$) & 7/7 \\
 & TT $-$ TG, log predictive density per visit & $+0.1060$ ($+0.0577$, $+0.1575$) & 7/7 \\
 & TT $-$ TG, class-label metric              & $+0.0379$ ($+0.0212$, $+0.0546$) & 7/7 \\
 & \multicolumn{3}{l}{\textbf{Verdict: supported.} Residual degrees of freedom $\approx 2.1$.} \\
\addlinespace[3pt]
\multicolumn{4}{@{}l}{\textit{(b) Are the patient slopes heavy-tailed? (Q2)}}\\
 & Original run: $\hat{R}$ / effective sample size & 1.018 / 248 & inconclusive \\
 & Reassessment: $\hat{R}$ / effective sample size & 1.0044 / 937 & converged \\
 & Full-cohort posterior median $\nu_{\mathrm{slope}}$ & 2.09 (1.48, 4.11) & --- \\
 & Per-split criterion met ($\nu$ upper bound $<10$) & 1 of 7 & threshold 5/7 \\
 & Null calibration, scenarios C1--C3 & 0 flags in 180 cohorts & --- \\
 & Positive control, scenario P & 55 of 60 detected & --- \\
 & \multicolumn{3}{l}{\textbf{Verdict: not supported.} No evidence of extreme slope tails.} \\
\addlinespace[3pt]
\multicolumn{4}{@{}l}{\textit{(c) Predictive value of each variant against the Gaussian model (Q3)}}\\
 & TG, class-label metric & $+0.0034$ ($+0.0007$, $+0.0063$) & 6/7 \\
 & GT, class-label metric & $+0.0382$ ($+0.0205$, $+0.0549$) & 7/7 \\
 & TT, class-label metric & $+0.0412$ ($+0.0229$, $+0.0586$) & 7/7 \\
\bottomrule
\end{tabular}
